\section{Introduction}
\label{sec:intro}

Understanding student engagement is central to effective teaching.
Disengagement---whether manifested as wandering attention, passive
participation, or behavioural withdrawal---predicts lower learning
outcomes and higher attrition~\cite{fredricks2004school}.
Traditionally, measuring engagement at scale requires trained human
observers~\cite{pianta2008classroom}, a resource-intensive process that
is difficult to replicate in large lecture halls or remote learning
environments.  The rapid proliferation of cameras in educational
settings, coupled with powerful visual recognition models, has therefore
sparked interest in automated engagement sensing~\cite{whitehill2014faces,
gupta2016daisee,thomas2018predicting,liao2021deep}.

Contemporary Vision-Language Models (VLMs) represent a potential leap
forward for this problem.  Unlike supervised classifiers that require
large labelled corpora specific to each school or cultural context,
models such as CLIP~\cite{radford2021learning},
BLIP~\cite{li2022blip}, GPT-4o~\cite{openai2024gpt4o}, and
LLaVA~\cite{liu2023visual} can be queried zero-shot via natural language,
with no task-specific fine-tuning.  Their ability to reason from open
prompts makes them attractive as \emph{plug-and-play} educational
observers.  Yet, to our knowledge, no prior work has systematically
benchmarked VLMs for classroom engagement recognition under controlled,
zero-shot conditions.

\begin{figure}[t]
    \centering
    \includegraphics[width=\linewidth]{figures/fig3_cross_dataset.pdf}
    \caption{Best zero-shot accuracy and Cohen's $\kappa$ per model on
    DAiSEE (individual student, video frames) versus SCB (classroom
    scene images).  Scene-level classification is substantially more
    tractable, with CLIP and GPT-4o achieving $\kappa\!\approx\!0.60$
    on SCB while remaining near-random on DAiSEE.}
    \label{fig:cross_dataset}
\end{figure}

\paragraph{Our contributions.}  This paper presents the following:

\begin{itemize}
  \item A zero-shot benchmark of four VLMs on two complementary
    datasets: the individual-student DAiSEE dataset~\cite{gupta2016daisee}
    and the scene-level Student Classroom Behaviour (SCB)
    dataset~\cite{wang2023scb}, using three systematically-designed
    prompt variants (minimal, rubric, chain-of-thought).

  \item A rigorous quantitative analysis documenting class collapse,
    near-random agreement on individual-level data ($\kappa \leq 0.10$),
    and substantial prompt sensitivity (accuracy variance up to 32 pp).

  \item The finding that scene-level engagement classification is
    significantly more tractable for VLMs, with CLIP achieving
    $\kappa = 0.60$ and GPT-4o achieving $\kappa = 0.58$ on SCB.

  \item Documentation of GPT-4o's safety-guardrail refusals blocking
    chain-of-thought engagement analysis for individual student faces,
    with practical implications for educational AI deployment.

  \item A self-consistency analysis showing that generative VLMs
    (GPT-4o, LLaVA) change their predictions for 27--32\% of images
    across repeated runs at temperature $T\!=\!0.7$, raising reliability
    concerns for classroom monitoring applications.
\end{itemize}

Our results paint a nuanced picture: while VLMs are not yet reliable
individual engagement observers, they demonstrate meaningful capability
at the scene level.  This gap between individual and scene-level
performance has direct implications for system design---aggregating
information spatially before querying a VLM may be a more productive
paradigm than per-student analysis.
